\section{Limitations and Conclusion}\label{sec:conclusion}

\textbf{Limitations.}
\begin{itemize}\setlength\itemsep{0pt}
\item \emph{Splicing is the main weakness.} Splices whose host and donor share the same processing history leave
no compression, noise or edge inconsistency; synthetic splices between two CASIA photographs were detected at
close to chance level (Section~\ref{sec:robustness}). Missed test-set forgeries were mostly splices
(Section~\ref{sec:test}).
\item \emph{Self-similar scenes cause false alarms.} Colonnades, lattices and repeated patterns mimic copy-moves.
All 13 confident false alarms on the test set had copy-move keypoint matches that survived RANSAC.
\item \emph{Localisation is coarse.} The fusion works on 8$\times$8 cells, and pixel F1 remains modest (0.198
under R; 0.333 for regions above 5\% of the image). Small regions ($<$1\%) are rarely found.
\item \emph{Limited external validation.} Only one external dataset, MICC-F220, could be obtained, and it contains
copy-moves from 11 scenes only. A cross-dataset splicing test is missing.
\item \emph{Remaining shortcuts.} The protocols remove most, not all, of the leak: under R the shortcut features
still reach AUC 0.712 on test, mainly through image size. We therefore report added value over the shortcuts, not
absolute accuracy alone.
\item \emph{Module limits.}
  \begin{itemize}\setlength\itemsep{0pt}
  \item The DCT periodicity test misses some double-quantisation patterns; it detected about 56\% of simulated
  double compressions.
  \item The keypoint matcher recovers at most one copied region per matching mode.
  \end{itemize}
\item \emph{Baseline scope.} The CNN comparison uses one architecture and one input representation, with less
tuning than the classical system.
\end{itemize}

\textbf{Conclusion.} CASIA~v2.0 as released does not measure what it is meant to measure: its labels can be
predicted from file metadata alone. Once that shortcut is removed and forensic features are required to beat the
cues that remain, an interpretable pipeline of classical DIP techniques still carries genuine tampering evidence:
\begin{itemize}\setlength\itemsep{0pt}
\item \emph{Classification.} Test AUC 0.805 under the strict protocol, +0.103 over the shortcut features.
\item \emph{Abstention.} The calibrated detector judges about half of the images, with 86\% accuracy.
\item \emph{Localisation.} It marks tampered regions and explains every decision with evidence maps and feature
attributions.
\item \emph{Transfer.} It reaches AUC 0.972 on an unseen copy-move dataset (MICC-F220, 11 forged scenes).
\end{itemize}
A standard ELA-CNN, by contrast, reaches its highest score (AUC 0.994) only on the leaking files, where shortcut
features alone reach 1.000. It trails the classical pipeline under both controlled protocols, and on that dataset
its AUC is consistent with chance.

The broader lesson concerns evaluation. Before a high score on a forensic benchmark is taken as evidence of
detection, the benchmark should be checked for shortcuts, and methods should be compared with what those shortcuts
alone achieve. Promising next steps are:
\begin{itemize}\setlength\itemsep{0pt}
\item cues aimed at same-history splices, such as sensor-noise consistency;
\item suppression of copy-move false alarms on repeated structure;
\item a cross-dataset splicing benchmark;
\item learned detectors trained and evaluated under the same leak controls.
\end{itemize}
